% GENERATED FILE. Edit canonical lessons, then run npm run generate:books.
% canonical-source-hash: fnv1a64:dd2d70ad8afe9022
% canonical-lessons: SA-C522-padanyasah, SA-C522-avartate, SA-C522-kinkini, SA-C522-anukaroti, SA-C522-karatalika

\chapter{A Song to Dance To}
\label{ch:sa-a-song-to-dance-to}
\hlchaptermodality{522}

\begin{chapteropening}
\textbf{By the end of this chapter:} \emph{I can talk about steps, turns and bells in a dance.}

Complete the last lesson of chapter 522: I can talk about steps, turns and bells in a dance.
\end{chapteropening}

\section[\texorpdfstring{\sk{पदन्यासः} (padanyāsaḥ)}{पदन्यासः (padanyāsaḥ)}]{\sk{पदन्यासः} (padanyāsaḥ) — a dance step}

\label{lesson:SA-C522-padanyasah}



\begin{quote}
Before the new one: say the Sanskrit for a tambourine, then the Sanskrit for a resonant sound.
\end{quote}

\subsection*{You'll want to know: \sk{पदन्यासः}}
\textbf{\sk{पदन्यासः}} — \emph{padanyāsaḥ} — \textquotedblleft{}a dance step\textquotedblright{}.

\begin{quote}
\sk{नर्तकस्य} \sk{पदन्यासः} \sk{सुन्दरः} \sk{अस्ति।}
\end{quote}

\begin{quote}
\emph{nartakasya padanyāsaḥ sundaraḥ asti.}
\end{quote}

\begin{quote}
The dancer's footwork is beautiful.
\end{quote}

A masculine noun.

\begin{cousinweb}
\textbf{\sk{पदन्यासः}} (\emph{padanyāsaḥ}) is built from \emph{pada}, \textquotedblleft{}a foot, a step\textquotedblright{}, and \emph{nyāsa}, \textquotedblleft{}a putting down\textquotedblright{}; MW gives \emph{padanyāsa} as putting down the feet, a step (MW).
\end{cousinweb}

\begin{grammarlens}[title={What you've built}]
One more word for describing what you lived and felt.
\end{grammarlens}

\subsection*{Your turn}
\begin{itemize}
\raggedright
  \item \emph{Say it:} \emph{padanyāsaḥ}
  \item \emph{Say it:} \emph{padanyāsaḥ}, once more
  \item \emph{Say it:} \emph{nartakasya padanyāsaḥ sundaraḥ asti.} — The dancer's footwork is beautiful.
  \item \emph{Recall:} say the Sanskrit for a tambourine, then the Sanskrit for a resonant sound, then say \emph{padanyāsaḥ} again
\end{itemize}

\begin{tcolorbox}[breakable,colback=teal!4,colframe=teal!35!black,title={Before you move on}]
What does \sk{पदन्यासः} mean? (\textquotedblleft{}A dance step\textquotedblright{}.) Say it once more.
\end{tcolorbox}
\section[\texorpdfstring{\sk{आवर्तते} (āvartate)}{आवर्तते (āvartate)}]{\sk{आवर्तते} (āvartate) — spins round; turns back}

\label{lesson:SA-C522-avartate}



\begin{quote}
Before the new one: say the Sanskrit for a resonant sound, then the Sanskrit for a dance step.
\end{quote}

\subsection*{You'll want to know: \sk{आवर्तते}}
\textbf{\sk{आवर्तते}} — \emph{āvartate} — \textquotedblleft{}spins round; turns back\textquotedblright{}.

\begin{quote}
\sk{नर्तकः} \sk{शीघ्रम्} \sk{आवर्तते।}
\end{quote}

\begin{quote}
\emph{nartakaḥ śīghram āvartate.}
\end{quote}

\begin{quote}
The dancer turns round quickly.
\end{quote}

Past (he …): \sk{सः} \sk{शीघ्रम्} \sk{आवृत्तवान्}, \emph{saḥ śīghram āvṛttavān}.

\begin{cousinweb}
\textbf{\sk{आवर्तते}} (\emph{āvartate}) is built on \emph{ā} and the root √vṛt, \textquotedblleft{}to turn\textquotedblright{}, which also gives \emph{vartate}; MW gives \emph{āvṛt} as to turn or draw round, to turn back (MW).
\end{cousinweb}

\begin{grammarlens}[title={What you've built}]
One more word for describing what you lived and felt.
\end{grammarlens}

\subsection*{Your turn}
\begin{itemize}
\raggedright
  \item \emph{Say it:} \emph{āvartate}
  \item \emph{Say it:} \emph{āvartate}, once more
  \item \emph{Say it:} \emph{nartakaḥ śīghram āvartate.} — The dancer turns round quickly.
  \item \emph{Recall:} say the Sanskrit for a resonant sound, then the Sanskrit for a dance step, then say \emph{āvartate} again
\end{itemize}

\begin{tcolorbox}[breakable,colback=teal!4,colframe=teal!35!black,title={Before you move on}]
What does \sk{आवर्तते} mean? (\textquotedblleft{}Spins round; turns back\textquotedblright{}.) Say it once more.
\end{tcolorbox}
\section[\texorpdfstring{\sk{किंकिणी} (kiṅkiṇī)}{किंकिणी (kiṅkiṇī)}]{\sk{किंकिणी} (kiṅkiṇī) — a small bell}

\label{lesson:SA-C522-kinkini}



\begin{quote}
Before the new one: say the Sanskrit for a dance step, then the Sanskrit for spins round.
\end{quote}

\subsection*{You'll want to know: \sk{किंकिणी}}
\textbf{\sk{किंकिणी}} — \emph{kiṅkiṇī} — \textquotedblleft{}a small bell\textquotedblright{}.

\begin{quote}
\sk{नर्तकस्य} \sk{पादे} \sk{किंकिण्यः} \sk{सन्ति।}
\end{quote}

\begin{quote}
\emph{nartakasya pāde kiṅkiṇyaḥ santi.}
\end{quote}

\begin{quote}
The dancer has small bells on his foot.
\end{quote}

A feminine noun.

\begin{cousinweb}
\textbf{\sk{किंकिणी}} (\emph{kiṅkiṇī}) is the word MW records for a small bell; MW gives no derivation, so its origin is uncertain (MW).
\end{cousinweb}

\begin{grammarlens}[title={What you've built}]
One more word for describing what you lived and felt.
\end{grammarlens}

\subsection*{Your turn}
\begin{itemize}
\raggedright
  \item \emph{Say it:} \emph{kiṅkiṇī}
  \item \emph{Say it:} \emph{kiṅkiṇī}, once more
  \item \emph{Say it:} \emph{nartakasya pāde kiṅkiṇyaḥ santi.} — The dancer has small bells on his foot.
  \item \emph{Recall:} say the Sanskrit for a dance step, then the Sanskrit for spins round, then say \emph{kiṅkiṇī} again
\end{itemize}

\begin{tcolorbox}[breakable,colback=teal!4,colframe=teal!35!black,title={Before you move on}]
What does \sk{किंकिणी} mean? (\textquotedblleft{}A small bell\textquotedblright{}.) Say it once more.
\end{tcolorbox}
\section[\texorpdfstring{\sk{अनुकरोति} (anukaroti)}{अनुकरोति (anukaroti)}]{\sk{अनुकरोति} (anukaroti) — imitates}

\label{lesson:SA-C522-anukaroti}



\begin{quote}
Before the new one: say the Sanskrit for spins round, then the Sanskrit for a small bell.
\end{quote}

\subsection*{You'll want to know: \sk{अनुकरोति}}
\textbf{\sk{अनुकरोति}} — \emph{anukaroti} — \textquotedblleft{}imitates\textquotedblright{}.

\begin{quote}
\sk{शिष्यः} \sk{गुरोः} \sk{पदन्यासम्} \sk{अनुकरोति।}
\end{quote}

\begin{quote}
\emph{śiṣyaḥ guroḥ padanyāsam anukaroti.}
\end{quote}

\begin{quote}
The pupil imitates the teacher's steps.
\end{quote}

Past (he …): \sk{सः} \sk{गुरुम्} \sk{अनुकृतवान्}, \emph{saḥ gurum anukṛtavān}.

\begin{cousinweb}
\textbf{\sk{अनुकरोति}} (\emph{anukaroti}) is built on \emph{anu}, \textquotedblleft{}after\textquotedblright{}, and the root √kṛ, \textquotedblleft{}to do\textquotedblright{}, which also gives \emph{karoti}: to do after another; MW gives \emph{anukṛ} as to imitate, to copy (MW).
\end{cousinweb}

\begin{grammarlens}[title={What you've built}]
One more word for describing what you lived and felt.
\end{grammarlens}

\subsection*{Your turn}
\begin{itemize}
\raggedright
  \item \emph{Say it:} \emph{anukaroti}
  \item \emph{Say it:} \emph{anukaroti}, once more
  \item \emph{Say it:} \emph{śiṣyaḥ guroḥ padanyāsam anukaroti.} — The pupil imitates the teacher's steps.
  \item \emph{Recall:} say the Sanskrit for spins round, then the Sanskrit for a small bell, then say \emph{anukaroti} again
\end{itemize}

\begin{tcolorbox}[breakable,colback=teal!4,colframe=teal!35!black,title={Before you move on}]
What does \sk{अनुकरोति} mean? (\textquotedblleft{}Imitates\textquotedblright{}.) Say it once more.
\end{tcolorbox}
\section[\texorpdfstring{\sk{करतालिका} (karatālikā)}{करतालिका (karatālikā)}]{\sk{करतालिका} (karatālikā) — clapping; applause}

\label{lesson:SA-C522-karatalika}



\begin{quote}
Before the new one: say the Sanskrit for a small bell, then the Sanskrit for imitates.
\end{quote}

\subsection*{You'll want to know: \sk{करतालिका}}
\textbf{\sk{करतालिका}} — \emph{karatālikā} — \textquotedblleft{}clapping; applause\textquotedblright{}.

\begin{quote}
\sk{नृत्यं} \sk{दृष्ट्वा} \sk{पिता} \sk{करतालिकां} \sk{करोति।}
\end{quote}

\begin{quote}
\emph{nṛtyaṃ dṛṣṭvā pitā karatālikāṃ karoti.}
\end{quote}

\begin{quote}
After the dance, father claps.
\end{quote}

A feminine noun.

\begin{cousinweb}
\textbf{\sk{करतालिका}} (\emph{karatālikā}) is built from \emph{kara}, \textquotedblleft{}the hand\textquotedblright{}, and \emph{tālikā}, \textquotedblleft{}clapping\textquotedblright{}, a word of the same family as \emph{tāla}, \textquotedblleft{}beating time\textquotedblright{}; MW gives \emph{karatālikā} as clapping the hands and beating time by clapping (MW).
\end{cousinweb}

\begin{grammarlens}[title={What you've built}]
One more word for describing what you lived and felt.
\end{grammarlens}

\subsection*{Your turn}
\begin{itemize}
\raggedright
  \item \emph{Say it:} \emph{karatālikā}
  \item \emph{Say it:} \emph{karatālikā}, once more
  \item \emph{Say it:} \emph{nṛtyaṃ dṛṣṭvā pitā karatālikāṃ karoti.} — After the dance, father claps.
  \item \emph{Recall:} say the Sanskrit for a small bell, then the Sanskrit for imitates, then say \emph{karatālikā} again
\end{itemize}

\begin{tcolorbox}[breakable,colback=teal!4,colframe=teal!35!black,title={Before you move on}]
What does \sk{करतालिका} mean? (\textquotedblleft{}Clapping; applause\textquotedblright{}.) Say it once more.
\end{tcolorbox}
